% ===== 第四部分：Fourier 分析 / 第1章：Fourier级数的基本性质 =====
% !TEX root = ../../main.tex
\setchapterimage{figure/P4/P4C1封面图.png}
\setchapterintro{本章从圆周上的 Fourier 系数与 Fourier 级数出发，建立正交性、唯一性以及函数光滑性与系数衰减之间的基本联系。随后引入卷积和好核，并借助 Fej\'er 核与 Poisson 核讨论 Ces\`aro 求和和 Abel 求和，从而为后续研究 Fourier 级数的收敛性奠定基础。}
\chapter{Fourier级数的基本性质}
% 第1节：主要的定义与一些例子
% !TEX root = ../../../main.tex
% 所属：第四部分：Fourier 分析 / 第1章：Fourier级数的基本性质
% 节：1.1 主要的定义与一些例子
% 正文按原文件顺序保存；由 ../chapter1.tex 调用。
\section{主要的定义与一些例子}
\begin{definition}
	如果 $f$ 是在给定长度为 $L$ 的区间 $[a,b]$ 上的函数,也就是说 $b-a=L$,那么 $f$ 的第 $n$ 个 Fourier 系数定义为
\[
\widehat f(n)=\frac1L\int_a^b f(x)e^{-2\pi i n x/L}\dd x,
\qquad n\in\Z .
\]
$f$ 的 Fourier 级数定义为
\[
\sum_{n=-\infty}^{\infty}\widehat f(n)e^{2\pi i n x/L}.
\]
我们有时把 $f$ 的 Fourier 系数记为 $a_n$,使用符号
\[
f(x)\sim \sum_{n=-\infty}^{\infty}a_n e^{2\pi i n x/L}
\]
表示右边的级数是 $f$ 的 Fourier 级数。

\end{definition}
例如,如果 $f$ 是区间 $[-\pi,\pi]$ 上的可积函数,那么 $f$ 的第 $n$ 个 Fourier 系数为
\[
\widehat f(n)=a_n
=\frac1{2\pi}\int_{-\pi}^{\pi}f(\theta)e^{-in\theta}\dd\theta,
\qquad n\in\Z,
\]
$f$ 的 Fourier 级数为
\[
f(\theta)
\sim \sum_{n=-\infty}^{\infty}a_n e^{in\theta}.
\]

在实际计算与应用中，我们经常采取Fourier展开的三角形式\[f(\theta)\sim \frac{a_0}{2}+\sum_{n=1}^{\infty}(a_n\cos nx+b_n\sin nx)\]

最后,考虑一个在区间 $[0,1]$ 上的函数 $g$,则
\[
\widehat g(n)=a_n
=\int_0^1 g(x)e^{-2\pi i n x}\dd x,
\qquad
g(x)\sim \sum_{n=-\infty}^{\infty}a_n e^{2\pi i n x}.
\]
其中,用 $x$ 作为从 $0$ 到 $1$ 的变量。

当然,如果 $f$ 最初定义在区间 $[0,2\pi]$ 上,那么令
\[
g(x)=f(2\pi x),\qquad x\in[0,1],
\]
由变量替换可知,$f$ 的第 $n$ 个 Fourier 系数和 $g$ 的第 $n$ 个 Fourier 系数是相同的。

出于方便考虑,在这一部分的笔记中,我们总是限制区间为$[-\pi,\pi]$或者$[0,2\pi]$


Fourier 级数是更大一类级数,也即三角级数的一部分。根据定义,三角级数可以表示为
\[
\sum_{n=-\infty}^{\infty}c_n e^{2\pi i n x/L},
\qquad c_n\in\C.
\]
如果三角级数只包含有限个非零项,即对于足够大的 $|n|$ 有 $c_n=0$,它就称作三角多项式；它的次数是 $|n|$ 的最大值,其中 $c_n\neq0$。

当 $N$ 是正整数时,$f$ 的 Fourier 级数的第 $N$ 个部分和定义为
\[
S_N(f)(x)=\sum_{n=-N}^{N}\widehat f(n)e^{2\pi i n x/L}.
\]
注意到通过定义,由于 $n$ 是从 $-N$ 到 $N$ 取值,故上述和是对称的。因此,可以将 Fourier 级数分解为正弦级数和余弦级数。因而,也可以将 Fourier 级数理解为当 $N$ 趋向无穷时这些对称和的极限。

一个需要注意到的事实是,任意一个三角多项式,在$S_N(f)(x)$的$\widehat f(\pm N)$非零的情况下,总可以由若干个$S_N(f)(x)$线性表出,这由线性代数中的知识显而易见


本章首先关心的问题是：\textbf{当 $N\to\infty$ 时,$S_N(f)$ 在何种意义下趋于 $f$。}

更具体地,为了便于讨论,假设函数 $f$ 总是 Riemann 可积的,并定义在 $[-\pi,\pi]$ 上。我们首先会问 $f$ 的 Fourier 级数的部分和是否逐点收敛于 $f$,也就是,对任意的 $\theta$,是否有

	\[\lim_{N\to\infty}S_N(f)(\theta)=f(\theta).\]

在正式的开始研究这个问题之前我们先给出一些例子


\begin{example}
	设
	\[
	f(\theta)=\theta,
	\qquad -\pi\leq\theta\leq\pi.
	\]
	计算 $f$ 的 Fourier 系数,并写出相应的 Fourier 级数。
\end{example}

\begin{solution}
	当 $n\neq0$ 时,
	\[
	\begin{aligned}
		\widehat f(n)
		&=\frac1{2\pi}\int_{-\pi}^{\pi}\theta e^{-in\theta}\dd\theta                                      \\
		&=\frac1{2\pi}\left(-\frac1{in}\right)
		\int_{-\pi}^{\pi}\theta\dd(e^{-in\theta})                                      \\
		&=-\frac1{2\pi in}
		\left(\left[\theta e^{-in\theta}\right]_{-\pi}^{\pi}
		-\int_{-\pi}^{\pi}e^{-in\theta}\dd\theta\right)                                  \\
		&=-\frac1{2\pi in}\cdot 2\pi(-1)^n
		=\frac{(-1)^{n+1}}{ni}.
	\end{aligned}
	\]
	当 $n=0$ 时,
	\[
	\widehat f(0)=\frac1{2\pi}\int_{-\pi}^{\pi}\theta\dd\theta=0.
	\]
	因此
	\[
	f(\theta)
	\sim \sum_{n\neq0}\frac{(-1)^{n+1}}{ni}e^{in\theta}.
	\]
	把正负指标合并,可得
	\[
	f(\theta)
	\sim 2\sum_{n=1}^{\infty}(-1)^{n+1}\frac{\sin n\theta}{n}.
	\]
\end{solution}

\begin{example}(Dirichlet 核)
	对 $x\in[-\pi,\pi]$,三角多项式
	\[
	D_N(x)=\sum_{n=-N}^{N}e^{inx}
	\]
	称作 $N$ 次 Dirichlet 核。证明它的闭形式为
	\[
	D_N(x)=\frac{\sin\left((N+\frac12)x\right)}{\sin(x/2)}.
	\]
\end{example}

\begin{solution}
	$D_N$ 在我们稍后谈到的理论中是至关重要的。注意到它的 Fourier 系数 $a_n$ 有性质
	\[
	a_n=
	\begin{cases}
		1,& |n|\leq N,\\
		0,& \text{其他}.
	\end{cases}
	\]
	令 $\omega=e^{ix}$。则
	\[
	D_N(x)=\sum_{n=-N}^{N}\omega^n
	=\frac{\omega^{-N}-\omega^{N+1}}{1-\omega}.
	\]
	于是
	\[
	\begin{aligned}
		D_N(x)
		&=\frac{\omega^{-N}-\omega^{N+1}}{1-\omega}                                      \\
		&=\frac{\omega^{-N-1/2}-\omega^{N+1/2}}{\omega^{-1/2}-\omega^{1/2}}                 \\
		&=\frac{\sin\left((N+\frac12)x\right)}{\sin(x/2)}.
	\end{aligned}
	\]
	这便给出了我们想要的结果。
\end{solution}

\begin{example}(Poisson 核)
	函数 $P_r(\theta)$ 称为 Poisson 核,是定义为 $\theta\in[-\pi,\pi]$ 和 $0\leq r<1$ 的绝对且一致收敛的级数
	\[
	P_r(\theta)=\sum_{n=-\infty}^{\infty}r^{|n|}e^{in\theta}.
	\]
	证明
	\[
	P_r(\theta)=\frac{1-r^2}{1-2r\cos\theta+r^2}.
	\]
\end{example}

\begin{solution}
	事实上,当 $\omega=re^{i\theta}$ 时,
	\[
	\begin{aligned}
		P_r(\theta)
		&=\sum_{n=0}^{\infty}\omega^n+
		\sum_{n=1}^{\infty}\overline{\omega}^{\,n}.
	\end{aligned}
	\]
	其中,这两个级数都绝对收敛。第一个和等于 $1/(1-\omega)$,同理,第二个等于 $\overline{\omega}/(1-\overline{\omega})$。合在一起,混合计算得到
	\[
	\begin{aligned}
		P_r(\theta)
		&=\frac1{1-\omega}+\frac{\overline{\omega}}{1-\overline{\omega}}                         \\
		&=\frac{1-\overline{\omega}+(1-\omega)\overline{\omega}}
		{(1-\omega)(1-\overline{\omega})}                                                         \\
		&=\frac{1-|\omega|^2}{|1-\omega|^2}                                                        \\
		&=\frac{1-r^2}{1-2r\cos\theta+r^2}.
	\end{aligned}
	\]
	这便得到了上述结果。Poisson 核会在今后函数的 Fourier 级数的 Abel 求和里再次出现。
\end{solution}

\begin{proposition}
	\begin{enumerate}
		\item 若$f$是偶函数,则$\widehat{f}(n)=\widehat{f}(-n)$
		\item 若$f$是奇函数,则$\widehat{f}(n)=-\widehat{f}(-n)$
	\end{enumerate}
\end{proposition}
\begin{proof}
	\begin{enumerate}
		\item \[\widehat{f}(n) =\frac{1}{2\pi}\int_{-\pi}^{\pi}f(\theta)e^{-in\theta}\dd \theta\]\[\widehat{f}(-n) =\frac{1}{2\pi}\int_{-\pi}^{\pi}f(\theta)e^{in\theta}\dd \theta=\frac{1}{2\pi}\int_{-\pi}^{\pi}f(-\theta)e^{-in\theta}\dd \theta=\widehat{f}(n)\]
		\item 同上
	\end{enumerate}
\end{proof}
若$f$是偶函数,可以写成
\[f(\theta)\sim \sum_{n=-\infty}^{\infty}a_n\cos n\theta\]
其中\[a_n =\frac{1}{2\pi}\int_{-\pi}^{\pi}f(\theta)\cos n\theta\dd \theta\]

若$f$是其函数,可以写成
\[f(\theta)\sim \sum_{n=-\infty}^{\infty}a_n\sin n\theta\]
其中\[a_n =\frac{1}{2\pi}\int_{-\pi}^{\pi}f(\theta)\sin
 n\theta\dd \theta\]





%
% 第2节：Fourier级数的唯一性
% !TEX root = ../../../main.tex
% 所属：第四部分：Fourier 分析 / 第1章：Fourier级数的基本性质
% 节：1.2 Fourier级数的唯一性
% 正文按原文件顺序保存；由 ../chapter1.tex 调用。
\section{Fourier级数的唯一性}
\begin{theorem}\label{Fourier级数的唯一性}
	假设 $f$ 是在圆周上可积的函数,且对所有 $n\in\Z$,有
	\[
	\widehat f(n)=0.
	\]
	若 $f$ 在点 $\theta_0$ 连续,则有
	\[
	f(\theta_0)=0.
	\]
\end{theorem}

\begin{proof}
	先假设 $f$ 是实值的。由于问题是周期的,并且可以通过平移把连续点移到原点,所以只需证明 $\theta_0=0$ 的情形。使用反证法,若 $f(0)\neq0$,不是一般性的可设$f(0)>0.$

	核心的思路是构造一个在零点取最大值的三角多项式列 $\{p_k\}$,使得当 $k\to\infty$ 时
	\[
	\int_{-\pi}^{\pi}p_k(\theta)f(\theta)\dd\theta\to\infty.
	\]
	但是根据 Fourier 系数全为零的假设,这些积分又都等于零,由此得到矛盾。
	
	由于 $f$ 在 $0$ 处连续,故可以选择 $0<\delta\leq \pi/2$,使得当 $|\theta|<\delta$ 时,
	\[
	f(\theta)>\frac{f(0)}2.
	\]
	令
	\[
	p(\theta)=\eps+\cos\theta,
	\]
	其中选择 $\eps>0$ 足够小,使当 $\delta\leq |\theta|\leq\pi$ 时,有
	\[
	|p(\theta)|<1-\frac\eps2.
	\]
	然后选择一个正数 $\eta$,且 $\eta<\delta$,使得对 $|\theta|<\eta$,有
	\[
	p(\theta)>1+\frac\eps2.
	\]
	最后令
	\[
	p_k(\theta)=\bigl[p(\theta)\bigr]^k.
	\]
	由$f\in \mathcal{R}(I)$,从而$f $有界,存在 $B>0$,使得对任意的 $\theta$,有 $|f(\theta)|\leq B$。每个 $p_k$ 都是三角多项式。由于对任意 $n$,有 $\widehat f(n)=0$,于是对任意 $k$,有
	\[
	\int_{-\pi}^{\pi}f(\theta)p_k(\theta)\dd\theta=0.
	\]
	另一方面,有估计
	\[
	\left|\int_{\delta\leq |\theta|\leq\pi}f(\theta)p_k(\theta)\dd\theta\right|
	\leq 2\pi B\left(1-\frac\eps2\right)^k.
	\]
	同时,选取 $\delta$ 使得当 $|\theta|<\delta$ 时,$p(\theta)$ 和 $f(\theta)$ 是非负的,于是
	\[
	\int_{\eta\leq |\theta|<\delta}f(\theta)p_k(\theta)\dd\theta\geq0.
	\]
	最后,在 $|\theta|<\eta$ 上,
	\[
	\begin{aligned}
		\int_{|\theta|<\eta}f(\theta)p_k(\theta)\dd\theta
		&\geq 2\eta\cdot \frac{f(0)}2\left(1+\frac\eps2\right)^k.
	\end{aligned}
	\]
	因此,当 $k\to\infty$ 时,
	\[
	\int_{-\pi}^{\pi}p_k(\theta)f(\theta)\dd\theta\to\infty,
	\]
	这与该积分恒等于零矛盾。所以实值情形下 $f(0)=0$。
	\begin{center}
		\begin{tikzpicture}[scale=1.0,>=Stealth]
			\draw[->] (-3.4,0) -- (3.4,0);
			\draw[->] (0,-0.2) -- (0,3.3);
			\draw[domain=-3.14:3.14,samples=200,smooth] plot (\x,{0.9+0.45*cos(\x r)});
			\draw[domain=-3.14:3.14,samples=300,smooth] plot (\x,{1.25*(0.1+cos(\x r))^3});
			\draw[domain=-1.2:1.2,samples=300,smooth] plot (\x,{2.45*(0.1+cos(\x r))^6});
			\node at (2.5,1.05) {$p$};
			\node at (0.95,1.9) {$p_3$};
			\node at (0.35,2.95) {$p_{6}$};
		\end{tikzpicture}
	\end{center}
	
	一般地,令
	\[
	f(\theta)=u(\theta)+iv(\theta),
	\]
	其中 $u$ 和 $v$ 是实值函数。若定义 $\overline f(\theta)=\overline{f(\theta)}$,那么有
	\[
	u(\theta)=\frac{f(\theta)+\overline f(\theta)}2,
	\qquad
	v(\theta)=\frac{f(\theta)-\overline f(\theta)}{2i}.
	\]
	又因为$\widehat{\overline f}(n)=\overline{\widehat f(-n)},$

	所以 $u$ 和 $v$ 的 Fourier 系数也全为零。由实值情形可得,在 $f$ 的连续点有
	\[
	u(\theta_0)=0,
	\qquad
	v(\theta_0)=0.
	\]
	因此 $f(\theta_0)=0$。
	
\end{proof}

根据上面的命题不难得到以下的命题
\begin{corollary}
	如果 $f$ 在圆周上连续且对所有 $n\in\Z$,有$\widehat f(n)=0,$则$f=0.$
	
\end{corollary}
\begin{corollary}
	如果 $f,g$ 在圆周上连续且对所有 $n\in\Z$,有$\widehat f(n)=\widehat g(n),$则$f=g.$
	
\end{corollary}
这也就是所谓的Fourier级数的唯一性

\begin{corollary}
	假设 $f$ 是在圆周上的连续函数且 $f$ 的 Fourier 系数绝对收敛,即
	\[
	\sum_{n=-\infty}^{\infty}|\widehat f(n)|<\infty.
	\]
	那么,$f$ 的 Fourier 级数一致收敛于 $f$,即
	\[
	\lim_{N\to\infty}S_N(f)(\theta)=f(\theta),
	\]
	且上述收敛关于 $\theta$ 是一致的。
\end{corollary}

\begin{proof}
由
	\[
	\sum_{n=-\infty}^{\infty}|\widehat f(n)|<\infty
	\]
	可知,$f$ 的 Fourier 系数的部分和绝对且一致收敛。事实上,由 Weierstrass $M$ 判别法,级数
	\[
	\sum_{n=-\infty}^{\infty}|\widehat f(n)e^{in\theta}| =\sum_{n=-\infty}^{\infty}|\widehat f(n)|<\infty
	\]
	在圆周上一致收敛。于是可以定义
	\[
	g(\theta)=\sum_{n=-\infty}^{\infty}\widehat f(n)e^{in\theta}
	=\lim_{N\to\infty}\sum_{n=-N}^{N}\widehat f(n)e^{in\theta}.
	\]
	由一致收敛和每一项的连续性可知,$g$ 在圆周上连续。
	
	进一步,由一致收敛可以交换无穷求和与积分,所以 $g$ 的 Fourier 系数恰好是 $\widehat f(n)$。具体地,
	\[
	\begin{aligned}
		\widehat g(n)
		&=\frac1{2\pi}\int_{-\pi}^{\pi}g(\theta)e^{-in\theta}\dd\theta                         \\
		&=\frac1{2\pi}\int_{-\pi}^{\pi}\sum_{m=-\infty}^{\infty}
		\widehat f(m)e^{im\theta}e^{-in\theta}\dd\theta                                             \\
		&=\frac1{2\pi}\sum_{m=-\infty}^{\infty}\widehat f(m)
		\int_{-\pi}^{\pi}e^{i(m-n)\theta}\dd\theta                                                   \\
		&=\widehat f(n).
	\end{aligned}
	\]
	因此 $f-g$ 的所有 Fourier 系数均为零。由上一推论,$f-g=0$,于是 $f=g$。也就是说,$f$ 的 Fourier 级数一致收敛于 $f$。
\end{proof}

\begin{corollary}\label{f(n)=O(1/n2)}
	假设 $f$ 是定义在圆周上的二次连续可微函数。当 $|n|\to\infty$ 时,有
	\[
	\widehat f(n)=O\left(\frac1{|n|^2}\right).
	\]
	则 $f$ 的 Fourier 级数绝对收敛且一致收敛于 $f$。
\end{corollary}

\begin{proof}
	当 $n\neq0$ 时,计算 Fourier 系数需要使用两次分部积分。取周期区间 $[0,2\pi]$,有
	\[
	\begin{aligned}
		2\pi\widehat f(n)
		&=\int_0^{2\pi}f(\theta)e^{-in\theta}\dd\theta                                      \\
		&=\left[f(\theta)\cdot \left(-\frac{e^{-in\theta}}{in}\right)\right]_0^{2\pi}
		+\frac1{in}\int_0^{2\pi}f'(\theta)e^{-in\theta}\dd\theta                              \\
		&=\frac1{in}\int_0^{2\pi}f'(\theta)e^{-in\theta}\dd\theta                              \\
		&=\frac1{in}\left[f'(\theta)\cdot\left(-\frac{e^{-in\theta}}{in}\right)\right]_0^{2\pi}
		+\frac1{(in)^2}\int_0^{2\pi}f''(\theta)e^{-in\theta}\dd\theta                          \\
		&=-\frac1{n^2}\int_0^{2\pi}f''(\theta)e^{-in\theta}\dd\theta.
	\end{aligned}
	\]
	由于 $f$ 和 $f'$ 是周期的,所以括号中的边界项均为零。因此
	\[
	2\pi |n|^2|\widehat f(n)|
	\leq \int_0^{2\pi}|f''(\theta)|\dd\theta
	\leq C,
	\]
	其中 $C$ 是与 $n$ 无关的常数。因为 $\sum 1/n^2$ 收敛,所以
	\[
	\sum_{n=-\infty}^{\infty}|\widehat f(n)|<\infty.
	\]
	由上一推论,$f$ 的 Fourier 级数绝对收敛且一致收敛于 $f$。
\end{proof}
进一步的,我们提高$f$的光滑性,得
\begin{corollary}\label{f(n)=O(1/n^k)}
	假设 $f\in C^k(\mathbb T)$,其中 $k\geq 1$。则当 $|n|\to\infty$ 时,
	\[
	\widehat f(n)=O\left(\frac1{|n|^k}\right).
	\]
	特别地,当 $k\geq 2$ 时,$f$ 的 Fourier 级数绝对收敛且一致收敛于 $f$。
\end{corollary}

我们还建立了以下重要等式：
\[
\widehat{f'}(n)=in\widehat f(n),
\qquad \forall n\in\Z.
\]
如果 $n\neq0$,这已经包含在上面的分部积分计算中；当 $n=0$ 时,
\[
\widehat{f'}(0)
=\frac1{2\pi}\int_0^{2\pi}f'(\theta)\dd\theta
=\frac{f(2\pi)-f(0)}{2\pi}=0=i0\widehat f(0),
\]
其中使用了 $f$ 的周期性。因此,如果 $f$ 是可微的且
\[
f\sim\sum a_ne^{in\theta},
\]
那么
\[
f'\sim\sum a_n in e^{in\theta}.
\]
同样,若 $f$ 是二次连续可微的,那么
\[
f''\sim\sum a_n(in)^2e^{in\theta}.
\]
在 $f$ 上附加更好的光滑性表现为 Fourier 系数衰减更迅速。
%
% 第3节：卷积
% !TEX root = ../../../main.tex
% 所属：第四部分：Fourier 分析 / 第1章：Fourier级数的基本性质
% 节：1.3 卷积
% 正文按原文件顺序保存；由 ../chapter1.tex 调用。
\section{卷积}
考虑\[
\begin{aligned}
	S_N(f)(x)
	&=\sum_{n=-N}^{N}\widehat f(n)e^{inx} \\
	&=\sum_{n=-N}^{N}\left(\frac1{2\pi}\int_{-\pi}^{\pi}f(y)e^{-iny}\dd y\right)e^{inx} \\
	&=\frac1{2\pi}\int_{-\pi}^{\pi}f(y)\left(\sum_{n=-N}^{N}e^{in(x-y)}\right)\dd y \\
\end{aligned}
\]
回忆$N$ 次 Dirichlet 核$D_N$ ,即
\[
D_N(x)=\sum_{n=-N}^{N}e^{inx}.
\]
可以给出以下的定义
\begin{definition}
	
给定两个在 $\R$ 上的、以 $2\pi$ 为周期的可积函数 $f$ 和 $g$,定义它们在 $[-\pi,\pi]$ 上的卷积为

	\[(f\ast g)(x)=\frac1{2\pi}\int_{-\pi}^{\pi}f(y)g(x-y)\dd y .\]

\end{definition}
因为两个可积函数的乘积也是可积的,故上面的积分对任意的 $x$ 都有意义。同样,由于函数具有周期性,通过改变变量可得
\[
(f\ast g)(x)=\frac1{2\pi}\int_{-\pi}^{\pi}f(x-y)g(y)\dd y .
\]

那么\[S_N(f)(x)=f\ast D_N(x)\]

因此,理解 $S_N(f)$ 的问题就简化为对卷积 $f\ast D_N$ 的理解。

\begin{proposition}
	假设 $f,g,h$ 是以 $2\pi$ 为周期的可积函数,则有：
	\[
	\begin{array}{ll}
		\mathrm{(i)}  & f\ast(g+h)=(f\ast g)+(f\ast h),\\[2mm]
		\mathrm{(ii)} & (cf)\ast g=c(f\ast g)=f\ast(cg),\quad \forall c\in\C,\\[2mm]
		\mathrm{(iii)}& f\ast g=g\ast f,\\[2mm]
		\mathrm{(iv)} & (f\ast g)\ast h=f\ast(g\ast h),\\[2mm]
		\mathrm{(v)}  & f\ast g \text{ 是连续的},\\[2mm]
		\mathrm{(vi)} & \widehat{f\ast g}(n)=\widehat f(n)\widehat g(n),\quad n\in\Z.
	\end{array}
	\]
\end{proposition}
前四点描述的是卷积的代数性质：线性性、交换性、结合性。性质 $\mathrm{(v)}$ 展示了一个重要原则：卷积具有平滑效果。这就是说,即使 $f$ 和 $g$ 只是 Riemann 可积的,$f\ast g$ 也是连续函数。最后 $\mathrm{(vi)}$ 是研究 Fourier 系数的关键,它说明卷积的 Fourier 系数等于 Fourier 系数的乘积。
\begin{proof}
	性质 $\mathrm{(i)}$ 和 $\mathrm{(ii)}$ 由积分的线性性立即得到。先假设 $f$ 和 $g$ 都是连续函数,则可以自由交换积分顺序(实际上在这一个问题中,仅仅有可积性,就可以交换积分次序)。对于性质 $\mathrm{(vi)}$,有
	\[
	\begin{aligned}
		\widehat{f\ast g}(n)
		&=\frac1{2\pi}\int_{-\pi}^{\pi}(f\ast g)(x)e^{-inx}\dd x \\
		&=\frac1{2\pi}\int_{-\pi}^{\pi}
		\left(\frac1{2\pi}\int_{-\pi}^{\pi}f(y)g(x-y)\dd y\right)e^{-inx}\dd x \\
		&=\frac1{2\pi}\int_{-\pi}^{\pi}f(y)e^{-iny}
		\left(\frac1{2\pi}\int_{-\pi}^{\pi}g(x-y)e^{-in(x-y)}\dd x\right)\dd y \\
		&=\frac1{2\pi}\int_{-\pi}^{\pi}f(y)e^{-iny}\widehat g(n)\dd y \\
		&=\widehat f(n)\widehat g(n).
	\end{aligned}
	\]
	
	为了证明性质 $\mathrm{(iii)}$,注意到如果 $F$ 是以 $2\pi$ 为周期的连续函数,则对任意 $x\in\R$,
	\[
	\int_{-\pi}^{\pi}F(y)\dd y
	=\int_{-\pi}^{\pi}F(x-y)\dd y.
	\]
	于是由变量替换可得
	\[
	(f\ast g)(x)=\frac1{2\pi}\int_{-\pi}^{\pi}f(x-y)g(y)\dd y=(g\ast f)(x).
	\]
	同样,性质 $\mathrm{(iv)}$ 在交换两个积分符号并作适当变量变换后得到。
	
	最后证明连续性。若 $f$ 和 $g$ 都连续,则 $g$ 在圆周上一致连续。任取 $\eps>0$,存在 $\delta>0$,使得当 $|s-t|<\delta$ 时,
	\[
	|g(s)-g(t)|<\eps.
	\]
	若 $|x_1-x_2|<\delta$,则对任意 $y$,有
	\[
	|(x_1-y)-(x_2-y)|<\delta.
	\]
	故
	\[
	\begin{aligned}
		|(f\ast g)(x_1)-(f\ast g)(x_2)|
		&\leq \frac1{2\pi}\int_{-\pi}^{\pi}|f(y)|\,|g(x_1-y)-g(x_2-y)|\dd y \\
		&\leq \frac{\eps}{2\pi}\int_{-\pi}^{\pi}|f(y)|\dd y.
	\end{aligned}
	\]
	于是 $f\ast g$ 是连续的。
	
	一般情况下,当 $f$ 和 $g$ 仅仅是可积函数时,需要用连续函数来逼近它们。\ref{C在R中一致有界稠密}给出这一点。

	对 $f$ 和 $g$ 应用\ref{C在R中一致有界稠密},得到近似的连续函数列 $\{f_k\}$ 和 $\{g_k\}$。于是
	\[
	f\ast g-f_k\ast g_k=(f-f_k)\ast g+f_k\ast(g-g_k).
	\]
	由引理的性质,当 $k\to\infty$ 时,
	\[
	\begin{aligned}
		|(f-f_k)\ast g(x)|
		&\leq \frac1{2\pi}\int_{-\pi}^{\pi}|f(x-y)-f_k(x-y)|\,|g(y)|\dd y \\
		&\leq \frac1{2\pi}\sup_y|g(y)|\int_{-\pi}^{\pi}|f(y)-f_k(y)|\dd y
		\to0.
	\end{aligned}
	\]
	所以 $(f-f_k)\ast g$ 关于 $x$ 一致收敛到 $0$。同理,$f_k\ast(g-g_k)$ 也一致收敛到 $0$。因此
	\[
	f_k\ast g_k\to f\ast g
	\]
	一致收敛。因为每个 $f_k\ast g_k$ 都是连续函数,故 $f\ast g$ 也是连续函数,从而得到了性质 $\mathrm{(v)}$。
	
	接下来证明性质 $\mathrm{(vi)}$。对每个固定的 $n$,由 $f_k\ast g_k$ 一致收敛于 $f\ast g$,得到
	\[
	\widehat{f_k\ast g_k}(n)\to\widehat{f\ast g}(n).
	\]
	另一方面,连续情形已经证明
	\[
	\widehat{f_k\ast g_k}(n)=\widehat f_k(n)\widehat g_k(n).
	\]
	并且
	\[
	|\widehat f(n)-\widehat f_k(n)|
	\leq \frac1{2\pi}\int_{-\pi}^{\pi}|f(x)-f_k(x)|\dd x\to0,
	\]
	类似地,$\widehat g_k(n)\to\widehat g(n)$。令 $k\to\infty$,便得到
	\[
	\widehat{f\ast g}(n)=\widehat f(n)\widehat g(n).
	\]
	最后,性质 $\mathrm{(iii)}$ 和性质 $\mathrm{(iv)}$ 也可以由连续函数情形再通过上述逼近过程得到。
\end{proof}

%
% 第4节：好核
% !TEX root = ../../../main.tex
% 所属：第四部分：Fourier 分析 / 第1章：Fourier级数的基本性质
% 节：1.4 好核
% 正文按原文件顺序保存；由 ../chapter1.tex 调用。
\section{好核}

卷积可以看成一种加权平均。若 $K_n$ 是以 $2\pi$ 为周期的函数,则
\[
(f\ast K_n)(x)=\frac1{2\pi}\int_{-\pi}^{\pi}f(x-y)K_n(y)\dd y.
\]
$K_n$起到了加权的作用,称做卷积核

\begin{definition}
	定义在圆周上的一列函数 $\{K_n(x)\}_{n=1}^{\infty}$ 满足以下性质：
	\[
	\begin{array}{ll}
		\mathrm{(a)} & \displaystyle \frac1{2\pi}\int_{-\pi}^{\pi}K_n(x)\dd x=1,\qquad n\geq1,\\[4mm]
		\mathrm{(b)} & \displaystyle \text{存在 }M>0,\text{ 对所有 }n\geq1,\text{ 有 }
		\int_{-\pi}^{\pi}|K_n(x)|\dd x\leq M,\\[4mm]
		\mathrm{(c)} & \displaystyle \text{对每一个 }\delta>0,\text{ 当 }n\to\infty\text{ 时,}
		\int_{\delta\leq |x|\leq\pi}|K_n(x)|\dd x\to0.
	\end{array}
	\]
	则称 $\{K_n(x)\}_{n=1}^{\infty}$ 为好核。
\end{definition}

\begin{theorem}
	假设 $\{K_n(x)\}_{n=1}^{\infty}$ 是好核列,且 $f$ 是圆周上的可积函数,则
	\[
	\lim_{n\to\infty}(f\ast K_n)(x)=f(x),
	\]
	其中 $f$ 在点 $x$ 连续。如果 $f$ 处处连续,那么上面的极限是一致收敛的。
\end{theorem}

\begin{proof}
	假设 $f$ 在点 $x$ 连续。给定 $\eps>0$,由连续性,选取 $\delta>0$,使得当 $|y|<\delta$ 时,有
	\[
	|f(x-y)-f(x)|<\eps.
	\]
	由好核的第一个性质,有
	\[
	\begin{aligned}
		(f\ast K_n)(x)-f(x)
		&=\frac1{2\pi}\int_{-\pi}^{\pi}K_n(y)f(x-y)\dd y-f(x) \\
		&=\frac1{2\pi}\int_{-\pi}^{\pi}K_n(y)[f(x-y)-f(x)]\dd y.
	\end{aligned}
	\]
	因此
	\[
	\begin{aligned}
		|(f\ast K_n)(x)-f(x)|
		&\leq \frac1{2\pi}\int_{|y|<\delta}|K_n(y)|\,|f(x-y)-f(x)|\dd y \\
		&\quad +\frac1{2\pi}\int_{\delta\leq |y|\leq\pi}|K_n(y)|\,|f(x-y)-f(x)|\dd y \\
		&\leq \frac{\eps}{2\pi}\int_{-\pi}^{\pi}|K_n(y)|\dd y
		+\frac{2B}{2\pi}\int_{\delta\leq |y|\leq\pi}|K_n(y)|\dd y,
	\end{aligned}
	\]
	其中 $B$ 是 $f$ 的界。由于好核的第二个性质,第一项被 $\eps M/(2\pi)$ 控制；由于好核的第三个性质,当 $n$ 足够大时,第二项可以任意小。因此存在常数 $C>0$,使得对充分大的 $n$,
	\[
	|(f\ast K_n)(x)-f(x)|<C\eps.
	\]
	这就证明了 $(f\ast K_n)(x)\to f(x)$。
	
	如果 $f$ 处处连续,则 $f$ 在圆周上一致连续,所以上面 $\delta$ 的选取可以独立于 $x$。于是上述估计对所有 $x$ 同时成立,从而得到一致收敛。
\end{proof}

%
% 第5节：Ces\`aro 和 Abel 求和：Fourier 级数的应用
% !TEX root = ../../../main.tex
% 所属：第四部分：Fourier 分析 / 第1章：Fourier级数的基本性质
% 节：1.5 Ces\`aro 和 Abel 求和：Fourier 级数的应用
% 正文按原文件顺序保存；由 ../chapter1.tex 调用。
\section{Ces\`aro 和 Abel 求和：Fourier 级数的应用}

因为 Fourier 级数可能在个别点处发散,为了克服这种缺点,我们寻求通过另一种方式来表示极限
\[
\lim_{N\to\infty}S_N(f)=f.
\]

假定给出一个复数级数
\[
c_0+c_1+c_2+\cdots=\sum_{k=0}^{\infty}c_k.
\]
定义第 $n$ 项和 $s_n$ 为
\[
s_n=\sum_{k=0}^{n}c_k.
\]
\begin{definition}
	我们定义前 $N$ 项和的平均值
\[
\sigma_N=\frac{s_0+s_1+\cdots+s_{N-1}}{N}
\]

量 $\sigma_N$ 被称为级数 $\{s_n\}$ 的 $N$ 次 Ces\`aro 平均,或级数 $\sum\limits_{k=0}^{\infty}c_k$ 的 $N$ 次 Ces\`aro 和。如果当 $N$ 趋于无穷时,$\sigma_N$ 收敛于一个极限 $s$,我们就说级数是 Ces\`aro 可和的,并称 $s$ 为它的 Ces\`aro 和。在函数级数的情形下,我们关心的是级数是逐点收敛的还是一致收敛的。

\end{definition}
我们早先提及 Dirichlet 核不属于好核的类。但它们的平均是性质很好的核。我们用对 Fourier 级数的 $N$ 次 Ces\`aro 和,具体定义如下：
\begin{definition}
	\[
\sigma_N(f)(x)=\frac{S_0(f)(x)+\cdots+S_{N-1}(f)(x)}{N}.
\]
因为 $S_N(f)=f\ast D_N$,故有
\[
\sigma_N(f)(x)=(f\ast F_N)(x),
\]
其中 $F_N(x)$ 是 $N$ 次 Fej\'er 核,定义如下：
\[
F_N(x)=\frac{D_0(x)+\cdots+D_{N-1}(x)}{N}=\sum_{\abs{k}\leqslant N}(1-\frac{|k|}{N})e^{ikx}.
\]
\end{definition}

\begin{lemma}
	可以得到
	\[
	F_N(x)=\frac1N\frac{\sin^2(Nx/2)}{\sin^2(x/2)},
	\]
	且 Fej\'er 核是好核。
\end{lemma}

\begin{proof}
	由定义,
	\[
	\begin{aligned}
		F_N(x)
		&=\frac1N\sum_{k=0}^{N-1}D_k(x)
		=\frac1N\sum_{k=0}^{N-1}\sum_{n=-k}^{k}e^{inx} \\
		&=\frac1N\sum_{|n|\leq N-1}(N-|n|)e^{inx}.
	\end{aligned}
	\]
	另一方面,
	\[
	\begin{aligned}
		\left|\sum_{j=0}^{N-1}e^{ijx}\right|^2
		&=\left(\sum_{j=0}^{N-1}e^{ijx}\right)
		\left(\sum_{k=0}^{N-1}e^{-ikx}\right) \\
		&=\sum_{|n|\leq N-1}(N-|n|)e^{inx}.
	\end{aligned}
	\]
	于是
	\[
	F_N(x)=\frac1N\left|\sum_{j=0}^{N-1}e^{ijx}\right|^2
	=\frac1N\left|\frac{1-e^{iNx}}{1-e^{ix}}\right|^2
	=\frac1N\frac{\sin^2(Nx/2)}{\sin^2(x/2)}.
	\]
	由上式可见 $F_N(x)\geq0$。又因为 $F_N$ 的第 $0$ 个 Fourier 系数为 $1$,所以
	\[
	\frac1{2\pi}\int_{-\pi}^{\pi}F_N(x)\dd x=1.
	\]
	因此
	\[
	\int_{-\pi}^{\pi}|F_N(x)|\dd x=\int_{-\pi}^{\pi}F_N(x)\dd x=2\pi.
	\]
	最后,给定 $\delta>0$,当 $\delta\leq |x|\leq\pi$ 时,$\sin^2(x/2)$ 有正的下界,于是
	\[
	0\leq F_N(x)\leq \frac{C_\delta}{N}.
	\]
	故
	\[
	\int_{\delta\leq |x|\leq\pi}F_N(x)\dd x\to0.
	\]
	所以 $\{F_N\}$ 是好核。
\end{proof}

\begin{theorem}
	如果 $f$ 在圆周上可积,那么 $f$ 的 Fourier 级数在 $f$ 的每个连续点是 Ces\`aro 和于 $f$ 的。并且,如果 $f$ 在圆周上连续,那么 $f$ 的 Fourier 级数是一致 Ces\`aro 和于 $f$ 的。
\end{theorem}

\begin{proof}
	由上一引理,$F_N$ 是好核。又因为
	\[
	\sigma_N(f)(x)=(f\ast F_N)(x),
	\]
	所以由好核逼近定理立即得到结论。
\end{proof}

\begin{corollary}
	若 $f$ 在圆周上可积且对所有 $n$,有 $\widehat f(n)=0$,则在 $f$ 的所有连续点都有
	\[
	f=0.
	\]
\end{corollary}

\begin{proof}
	这显然是成立的。因为 $f$ 的所有部分和都是零,所以其 Ces\`aro 平均也是零。由上一定理,在 $f$ 的连续点处,Ces\`aro 平均收敛到 $f$,于是 $f=0$。
\end{proof}

\begin{corollary}
	圆周上的连续函数可以被三角多项式一致逼近。更精确地说,给定任意 $\eps>0$,存在三角多项式 $P$,使得对任意的 $-\pi\leq x\leq\pi$,
	\[
	|f(x)-P(x)|<\eps
	\]
	成立。
\end{corollary}

\begin{proof}
	由于 Fourier 部分和是三角多项式,所以 Ces\`aro 平均 $\sigma_N(f)$ 也是三角多项式,从而直接由上一定理的一致收敛性得到结论。这个结论就是周期函数的 Weierstrass 逼近定理。
\end{proof}

\begin{definition}
	如果对任意的 $0<r<1$,级数
\[
A(r)=\sum_{k=0}^{\infty}c_k r^k
\]
收敛,且
\[
\lim_{r\to1}A(r)=s,
\]
则复数 $s$ 称为级数 $\sum\limits_{k=0}^{\infty}c_k$ 的 Abel 求和,$A(r)$ 称为此级数的 Abel 平均。如果这个级数收敛于 $s$,那么它是 Abel 和于 $s$
\end{definition}
 另外,Abel 和方法甚至比 Ces\`aro 方法更强。也就是说,当级数是 Ces\`aro 和时,它也是 Abel 和于相同的数。然而,如果考虑级数
\[
1-2+3-4+5-\cdots=\sum_{k=0}^{\infty}(-1)^k(k+1),
\]
由于
\[
A(r)=\sum_{k=0}^{\infty}(-1)^k(k+1)r^k=\frac1{(1+r)^2},
\]
故可以证明它是 Abel 和于 $1/4$,但这个级数不是 Ces\`aro 和的。
\begin{proposition}
	Abel 求和要比一般求和以及 Ces\`aro 求和更强,具体如下
	
	\begin{enumerate}
		\item[(a)] 若复数级数$\sum\limits_{n=1}^{\infty}c_n$
		收敛于一个有限数 $s$,则它的 Abel 和为 $s$
		\item[(b)] 存在 Abel 可求和但是不收敛的级数。
		\item[(c)] 类似地,证明：若$\sum\limits_{n=1}^{\infty}c_n$的 Ces\`aro 和为 $\sigma$,则它的 Abel 和也为 $\sigma$。
		\item[(d)] 给出一个有 Abel 和,但是 Ces\`aro 和发散的例子。
	\end{enumerate}
	这些结果可以总结为
	\[
	\text{收敛}
	\Longrightarrow
	\text{Ces\`aro 可求和}
	\Longrightarrow
	\text{Abel 可求和},
	\]
	其中的箭头均无法逆向成立。
\end{proposition}

\begin{solution}
	\textbf{(a)}
	
	
	先假设$ \sum\limits_{n=1}^{\infty}c_n=0.$令$s_N=\sum\limits_{n=1}^{N}c_n,$则 $s_N\to0$,从而数列 $\{s_N\}$ 有界。
	
	注意到
	\[
	c_n=s_n-s_{n-1},
	\qquad s_0=0,
	\]
	所以
	\[\sum_{n=1}^{N}c_nr^n
	=\sum_{n=1}^{N}(s_n-s_{n-1})r^n
	=\sum_{n=1}^{N}s_nr^n
	-r\sum_{n=1}^{N-1}s_nr^n
	=(1-r)\sum_{n=1}^{N}s_nr^n+s_Nr^{N+1}
	\]
	
	固定 $0<r<1$。由于 $\{s_N\}$ 有界,
	\[
	s_Nr^{N+1}\to0.
	\]
	令 $N\to\infty$,得到
	\[
	\sum_{n=1}^{\infty}c_nr^n
	=(1-r)\sum_{n=1}^{\infty}s_nr^n.
	\]
	
	下面证明右端在 $r\to1^-$ 时趋于 $0$。给定 $\varepsilon>0$,由 $s_n\to0$,存在正整数 $N_0$,使得当 $n>N_0$ 时,
	\[
	|s_n|<\varepsilon.
	\]
	于是
	\[
	\begin{aligned}
		(1-r)\left|\sum_{n=1}^{\infty}s_nr^n\right|
		&\leq
		(1-r)\sum_{n=1}^{N_0}|s_n|r^n
		+(1-r)\sum_{n=N_0+1}^{\infty}|s_n|r^n\\
		&\leq
		(1-r)\sum_{n=1}^{N_0}|s_n|r^n
		+\varepsilon(1-r)\sum_{n=N_0+1}^{\infty}r^n.
	\end{aligned}
	\]
	当 $r\to1^-$ 时,第一项趋于 $0$,而第二项满足
	\[
	\varepsilon(1-r)\sum_{n=N_0+1}^{\infty}r^n
	=\varepsilon r^{N_0+1}
	\leq\varepsilon.
	\]
	因此
	\[
	\lim_{r\to1^-}(1-r)\sum_{n=1}^{\infty}s_nr^n=0.
	\]
	所以
	\[
	\lim_{r\to1^-}\sum_{n=1}^{\infty}c_nr^n=0.
	\]
	
	综上,若
	\[
	\sum_{n=1}^{\infty}c_n=s,
	\]
	则该级数的 Abel 和也是 $s$。
	再说明为什么只需考虑 $s=0$ 的情形。若
	\[
	\sum_{n=1}^{\infty}c_n=s,
	\]
	定义
	\[
	d_1=c_1-s,\qquad d_n=c_n\quad(n\geq2).
	\]
	则
	\[
	\sum_{n=1}^{\infty}d_n=0,
	\]
	并且
	\[
	\sum_{n=1}^{\infty}c_nr^n
	=\sum_{n=1}^{\infty}d_nr^n+sr.
	\]
	因此,只要证明和为 $0$ 的收敛级数的 Abel 和为 $0$,便有
	\[
	\lim_{r\to1^-}\sum_{n=1}^{\infty}c_nr^n
	=0+s=s.
	\]
	
	\textbf{(b)}取	$c_n=(-1)^n.$相应的级数为$-1+1-1+1-\cdots.$其部分和满足
	\[
	s_{2m}=0,\qquad s_{2m-1}=-1,
	\]
	因此该级数不收敛。
	
	但是,当 $0<r<1$ 时,
	\[	\sum_{n=1}^{\infty}c_nr^n
	=\sum_{n=1}^{\infty}(-1)^nr^n
	=-r+r^2-r^3+\cdots
	=-\frac{r}{1+r}.
	\]
	所以
	\[
	\lim_{r\to1^-}\sum_{n=1}^{\infty}c_nr^n
	=-\frac12.
	\]
	因此该级数虽然不收敛,却是 Abel 可求和的,且 Abel 和为$-\frac12.$
	
	事实上,该级数也是 Ces\`aro 可求和的。因为
	\[
	\sigma_N=\frac{s_1+\cdots+s_N}{N},
	\]
	所以
	\[
	\sigma_{2m}
	=\frac{-m}{2m}
	=-\frac12,
	\]
	而
	\[
	\sigma_{2m-1}
	=\frac{-m}{2m-1}
	\to-\frac12.
	\]
	故其 Ces\`aro 和也为 $-\dfrac12$。
	
	\textbf{(c)}设
	\[
	s_n=c_1+\cdots+c_n,
	\qquad
	\sigma_n=\frac{s_1+\cdots+s_n}{n},
	\]
	并假设$\sigma_n\to\sigma.$
	
	类似于第 (a) 问,只需考虑 $\sigma=0$ 的情形。事实上,若定义
	\[
	d_1=c_1-\sigma,\qquad d_n=c_n\quad(n\geq2),
	\]
	则新级数的部分和为$t_n=s_n-\sigma,$其 Ces\`aro 平均为
	\[
	\frac{t_1+\cdots+t_n}{n}
	=\sigma_n-\sigma\to0.
	\]
	同时
	\[
	\sum_{n=1}^{\infty}c_nr^n
	=\sum_{n=1}^{\infty}d_nr^n+\sigma r.
	\]
	因此只需证明 Ces\`aro 和为 $0$ 时,Abel 和也为 $0$。
	
	下面假设$\sigma_n\to0.$由$n\sigma_n=s_1+\cdots+s_n$可得
	\[
	s_n=n\sigma_n-(n-1)\sigma_{n-1},
	\]
	其中规定 $\sigma_0=0$。
	
	于是
	\[\sum_{n=1}^{\infty}s_nr^n
	=\sum_{n=1}^{\infty}
	\bigl(n\sigma_n-(n-1)\sigma_{n-1}\bigr)r^n
	=\sum_{n=1}^{\infty}n\sigma_nr^n
	-r\sum_{n=1}^{\infty}n\sigma_nr^n
	=(1-r)\sum_{n=1}^{\infty}n\sigma_nr^n.
	\]
	由第 (a) 问中的分部求和公式,
	\[
	\sum_{n=1}^{\infty}c_nr^n
	=(1-r)\sum_{n=1}^{\infty}s_nr^n.
	\]
	因此
	\[
	\sum_{n=1}^{\infty}c_nr^n
	=(1-r)^2\sum_{n=1}^{\infty}n\sigma_nr^n.
	\]
	
	给定 $\varepsilon>0$,由 $\sigma_n\to0$,存在正整数 $N_0$,使得当 $n>N_0$ 时,
	\[
	|\sigma_n|<\varepsilon.
	\]
	于是
	\[(1-r)^2
	\left|\sum_{n=1}^{\infty}n\sigma_nr^n\right|
	\leq
	(1-r)^2\sum_{n=1}^{N_0}n|\sigma_n|r^n
	+
	\varepsilon(1-r)^2
	\sum_{n=N_0+1}^{\infty}nr^n.\]
	当 $r\to1^-$ 时,第一项趋于 $0$。又因为
	\[
	\sum_{n=1}^{\infty}nr^n=\frac{r}{(1-r)^2},
	\]
	所以第二项不超过
	\[
	\varepsilon(1-r)^2
	\sum_{n=1}^{\infty}nr^n
	=\varepsilon r
	\leq\varepsilon.
	\]
	故
	\[
	\lim_{r\to1^-}
	(1-r)^2\sum_{n=1}^{\infty}n\sigma_nr^n=0.
	\]
	
	因此,当 Ces\`aro 和为 $0$ 时,Abel 和也为 $0$。由前面的证明可知,若级数的 Ces\`aro 和为 $\sigma$,则其 Abel 和也是 $\sigma$。
	
	\textbf{(d)}取$c_n=(-1)^{n-1}n,$即考虑级数$	1-2+3-4+5-\cdots.$当 $0<r<1$ 时,由
	\[
	\sum_{n=0}^{\infty}(-r)^n=\frac{1}{1+r}
	\]
	两边对 $r$ 求导,得到
	\[
	\sum_{n=1}^{\infty}(-1)^{n-1}nr^{n-1}
	=\frac{1}{(1+r)^2}.
	\]
	两边乘以 $r$,得到
	\[
	\sum_{n=1}^{\infty}(-1)^{n-1}nr^n
	=\frac{r}{(1+r)^2}.
	\]
	因此
	\[
	\lim_{r\to1^-}
	\sum_{n=1}^{\infty}(-1)^{n-1}nr^n
	=\frac14.
	\]
	所以该级数是 Abel 可求和的,且 Abel 和为$\dfrac14.$
	
	
	下面证明它不是 Ces\`aro 可求和的。首先证明：若级数 $\sum c_n$ 是 Ces\`aro 可求和的,则$\dfrac{c_n}{n}\to0.$
	若 $\sigma_n\to\sigma$,则
	\[
	s_n=n\sigma_n-(n-1)\sigma_{n-1},
	\]
	从而
	\[
	\frac{s_n}{n}
	=\sigma_n-\frac{n-1}{n}\sigma_{n-1}
	\to\sigma-\sigma=0.
	\]
	又因为
	\[
	c_n=s_n-s_{n-1},
	\]
	所以
	\[
	\frac{c_n}{n}
	=\frac{s_n}{n}
	-\frac{n-1}{n}\frac{s_{n-1}}{n-1}
	\to0.
	\]
	
	但是对于$c_n=(-1)^{n-1}n,$有	$\dfrac{c_n}{n}=(-1)^{n-1},$
	
	它不趋于 $0$。因此级数$1-2+3-4+5-\cdots$不是 Ces\`aro 可求和的。
	
	综上,
	\[
	\boxed{
		\text{收敛}
		\Longrightarrow
		\text{Ces\`aro 可求和}
		\Longrightarrow
		\text{Abel 可求和}
	}
	\]
	且两个箭头均不能反向成立。
\end{solution}
下面介绍Tauber的一个定理,对上面的$c_n$
,添加一个条件后,上面的箭头就是可逆的
\begin{proposition}
	若$\sum c_n$的Ces\'aro和为$\sigma$,且$c_n = o(\dfrac{1}{n})$,则$\sum c_n$收敛与$\sigma$,对Abel和,有同样的结论
\end{proposition}
\begin{proof}
	考虑\[s_n- \sigma_n =\frac{(n-1)c_n +\cdots+c_2}{n}\]
	由$c_n = o(\dfrac{1}{n})$,$\forall \eps >0,\exists N_0,\text{使得}\forall n>N_0,c_n<\dfrac{\eps}{n}$
	那么
	\[s_n- \sigma_n =\frac{(n-1)c_n +\cdots+c_2}{n} =\frac{c_2+\cdots+(N_0-1)c_{N_0}}{n}+\frac{N_0 c_{N_0+1}\cdots+(n-1)c_n}{n}\]
	对第一项,存在$N>N_0$使,$\forall n>N$,$ \dfrac{c_2+\cdots+(N_0-1)c_{N_0}}{n}<\dfrac{\eps}{2}$对第二项
	\[\frac{N_0 c_{N_0+1}\cdots+(n-1)c_n}{n}<(1-\frac{N_0}{n})\frac{\eps}{2}\],因此
	\[s_n- \sigma_n< c\eps\]
	从而\[\lim\limits_{n\to \infty}s_n- \sigma_n =0\]
	对Abel和,类似的,我们考虑\[\sum_{n=1}^{N}c_n-\sum_{n=1}^{N}c_n r^n\]并且不妨取$r=1-\dfrac{1}{N}$
	\[\sum_{n=1}^{N}c_n-\sum_{n=1}^{N}c_n r^n<\sum_{n=1}^{N}c_n\frac{n}{N}\]类似于上面的方法,不难得出\[\sum_{n=1}^{N}c_n-\sum_{n=1}^{N}c_n r^n\to0\]
\end{proof}


\begin{definition}
	定义函数
\[
f(\theta)\sim\sum_{n=-\infty}^{\infty}a_ne^{in\theta}
\]
的 Abel 平均为
\[
A_r(f)(\theta)=\sum_{n=-\infty}^{\infty}r^{|n|}a_ne^{in\theta}.
\]
\end{definition}
因为指数 $n$ 取正值和负值,很自然地出现 $r^{|n|}$。并且当 $r>0$ 时,$c_n=a_ne^{in\theta}+a_{-n}e^{-in\theta}$,所以 Fourier 级数的 Abel 平均对应于普通级数的 Abel 平均的定义。

因为 $f$ 是可积的,$|a_n|$ 一致有界,所以对每个 $0\leq r<1$,$A_r(f)$ 绝对收敛且一致收敛。正如 Ces\`aro 平均的情形,关键是这些 Abel 平均可以写成卷积的形式,即
\[
A_r(f)(\theta)=(f\ast P_r)(\theta),
\]
其中 $P_r(\theta)$ 是 Poisson 核,定义如下：

	$$P_r(\theta)=\sum_{n=-\infty}^{\infty}r^{|n|}e^{in\theta}.$$

事实上,有
\[
\begin{aligned}
	A_r(f)(\theta)
	&=\sum_{n=-\infty}^{\infty}r^{|n|}a_ne^{in\theta} \\
	&=\sum_{n=-\infty}^{\infty}r^{|n|}
	\left(\frac1{2\pi}\int_{-\pi}^{\pi}f(\varphi)e^{-in\varphi}\dd\varphi\right)e^{in\theta} \\
	&=\frac1{2\pi}\int_{-\pi}^{\pi}f(\varphi)
	\left(\sum_{n=-\infty}^{\infty}r^{|n|}e^{in(\theta-\varphi)}\right)\dd\varphi,
\end{aligned}
\]
其中,交换积分和求和顺序是因为该级数一致收敛。

\begin{lemma}
	如果 $0\leq r<1$,那么有
	\[
	P_r(\theta)=\frac{1-r^2}{1-2r\cos\theta+r^2}.
	\]
	当 $r\to1$ 时,Poisson 核是好核。
\end{lemma}

\begin{proof}
	等式
	\[
	P_r(\theta)=\frac{1-r^2}{1-2r\cos\theta+r^2}
	\]
	已经在前面关于 Poisson 核的例子中给出。注意到由这个表达式可知 $P_r(\theta)>0$。此外,由于 $P_r$ 的第 $0$ 个 Fourier 系数为 $1$,得到
	\[
	\frac1{2\pi}\int_{-\pi}^{\pi}P_r(\theta)\dd\theta=1.
	\]
	因此
	\[
	\int_{-\pi}^{\pi}|P_r(\theta)|\dd\theta=2\pi.
	\]
	最后,给定 $\delta>0$。当 $1/2\leq r\leq1$ 且 $\delta\leq |\theta|\leq\pi$ 时,存在常数 $c_\delta>0$,使得
	\[
	1-2r\cos\theta+r^2\geq c_\delta.
	\]
	于是
	\[
	0\leq P_r(\theta)\leq \frac{1-r^2}{c_\delta}.
	\]
	当 $r\to1$ 时,上式右边趋于 $0$,故
	\[
	\int_{\delta\leq |\theta|\leq\pi}P_r(\theta)\dd\theta\to0.
	\]
	所以当 $r\to1$ 时,Poisson 核构成好核。
\end{proof}

\begin{theorem}
	在圆周上可积函数的 Fourier 级数在每个连续点 Abel 和于 $f$。此外,如果 $f$ 在圆周上连续,那么 $f$ 的 Fourier 级数一致 Abel 和于 $f$。
\end{theorem}

\begin{proof}
	由上一引理,$P_r$ 在 $r\to1$ 时构成好核。又因为
	\[
	A_r(f)(\theta)=(f\ast P_r)(\theta),
	\]
	所以由好核逼近定理可得结论。
\end{proof}
下面,我们考虑$f$在$\theta$处不连续的情况,我们只考虑跳跃点,即第一类间断点
\begin{proposition}

	\begin{enumerate}
		\item[(a)] 若 $f$ 在 $\theta$ 处是跃点,则
		\[
		\lim_{r\to1^-}A_r(f)(\theta)
		=\frac{f(\theta^+)+f(\theta^-)}{2}.
		\]
		
		\item[(b)] 类似的：若 $f$ 在 $\theta$ 处有跳跃间断,
		则 $f$ 在 $\theta$ 处的 Fourier 级数的 Ces\`aro 和为
		\[
		\frac{f(\theta^+)+f(\theta^-)}{2}.
		\]
	\end{enumerate}
	
\end{proposition}


\begin{proof}
	记
	\[
	\alpha=f(\theta^+),\qquad
	\beta=f(\theta^-),\qquad
	L=\frac{\alpha+\beta}{2}.
	\]
	
	\textbf{(a)}
	\[
	A_r(f)(\theta)
	=(P_r*f)(\theta)
	=\frac{1}{2\pi}\int_{-\pi}^{\pi}
	f(\theta-t)P_r(t)\dd t,
	\]
	其中 Poisson 核为
	\[
	P_r(t)=\frac{1-r^2}{1-2r\cos t+r^2}.
	\]
	
	由于 $P_r$ 是偶函数,且
	\[
	\frac{1}{2\pi}\int_{-\pi}^{\pi}P_r(t)\dd t=1,
	\]
	所以
	\[
	\frac{1}{2\pi}\int_{-\pi}^{0}P_r(t)\dd t
	=
	\frac{1}{2\pi}\int_{0}^{\pi}P_r(t)\dd t
	=\frac12.
	\]
	
	在积分区间 $[-\pi,0]$ 上作变量代换 $t\mapsto-t$,并利用
	$P_r(-t)=P_r(t)$,可得
	\[
	\begin{aligned}
		A_r(f)(\theta)
		&=\frac{1}{2\pi}\int_0^\pi
		f(\theta-t)P_r(t)\dd t
		+\frac{1}{2\pi}\int_0^\pi
		f(\theta+t)P_r(t)\dd t\\
		&=\frac{1}{2\pi}\int_0^\pi
		\bigl(f(\theta-t)+f(\theta+t)\bigr)P_r(t)\dd t.
	\end{aligned}
	\]
	
	另一方面,由
	\[
	\frac{1}{2\pi}\int_0^\pi P_r(t)\dd t=\frac12,
	\]
	有
	\[
	L
	=\frac{\alpha+\beta}{2}
	=\frac{1}{2\pi}\int_0^\pi
	(\alpha+\beta)P_r(t)\dd t.
	\]
	因此
	\[
	\begin{aligned}
		A_r(f)(\theta)-L
		&=\frac{1}{2\pi}\int_0^\pi
		\bigl[
		f(\theta+t)-\alpha
		+f(\theta-t)-\beta
		\bigr]P_r(t)\dd t.
	\end{aligned}
	\]
	
	给定 $\varepsilon>0$。由$\lim\limits_{\substack{t\to0\\t>0}}f(\theta+t)=\alpha,
	\lim\limits_{\substack{t\to0\\t>0}}f(\theta-t)=\beta,$
	存在 $\delta>0$,使得当 $0<t<\delta$ 时,
	\[
	|f(\theta+t)-\alpha|<\varepsilon,
	\qquad
	|f(\theta-t)-\beta|<\varepsilon.
	\]
	
	将积分分为 $(0,\delta)$ 和 $(\delta,\pi)$ 两部分。对于第一部分,
	由 $P_r(t)\geq0$,有
	\[\frac{1}{2\pi}\int_0^\delta
		\left|
		f(\theta+t)-\alpha
		+f(\theta-t)-\beta
		\right|P_r(t)\dd t\leq
		\frac{2\varepsilon}{2\pi}
		\int_0^\delta P_r(t)\dd t\leq
		\frac{2\varepsilon}{2\pi}
		\int_0^\pi P_r(t)\dd t
		=\varepsilon.
	\]
	
	对于第二部分,令
	\[
	C_\delta
	=
	\int_\delta^\pi
	\bigl(
	|f(\theta+t)|+|f(\theta-t)|+|\alpha|+|\beta|
	\bigr)\dd t.
	\]
	由于 $f$ 可积,故 $C_\delta<\infty$。于是
	\[\frac{1}{2\pi}\int_\delta^\pi
		\left|
		f(\theta+t)-\alpha
		+f(\theta-t)-\beta
		\right|P_r(t)\dd t\leq
		\frac{C_\delta}{2\pi}
		\sup_{\delta\leq t\leq\pi}P_r(t).\]
	
	当 $\delta\leq t\leq\pi$ 时,
	\[
	1-2r\cos t+r^2
	=(1-r)^2+2r(1-\cos t)
	\geq 2r(1-\cos\delta).
	\]
	因此,当 $r$ 充分接近 $1$ 时,
	\[
	P_r(t)
	\leq
	\frac{1-r^2}{2r(1-\cos\delta)},
	\qquad \delta\leq t\leq\pi.
	\]
	从而
	\[
	\sup_{\delta\leq t\leq\pi}P_r(t)\to0,
	\qquad r\to1^-.
	\]
	
	由此可得
	\[
	\limsup_{r\to1^-}|A_r(f)(\theta)-L|
	\leq\varepsilon.
	\]
	由于 $\varepsilon>0$ 任意,所以
	\[
	\lim_{r\to1^-}A_r(f)(\theta)
	=L
	=\frac{f(\theta^+)+f(\theta^-)}{2}.
	\]

	
	\textbf{(b)}
	由 Fej\'er 平均的卷积表示,
	\[
	\sigma_N(f)(\theta)
	=(F_N*f)(\theta)
	=\frac{1}{2\pi}\int_{-\pi}^{\pi}
	f(\theta-t)F_N(t)\dd t,
	\]
	其中 Fej\'er 核为
	\[
	F_N(t)
	=\frac{1}{N}
	\left(
	\frac{\sin(Nt/2)}{\sin(t/2)}
	\right)^2.
	\]
	
	Fej\'er 核具有下列性质：
	\[
	F_N(t)\geq0,\qquad F_N(-t)=F_N(t),
	\]
	并且
	\[
	\frac{1}{2\pi}\int_{-\pi}^{\pi}F_N(t)\dd t=1.
	\]
	因此
	\[
	\frac{1}{2\pi}\int_0^\pi F_N(t)\dd t=\frac12.
	\]
	
	与第 (a) 问相同,在负半区间作变量代换,得到
	\[
	\sigma_N(f)(\theta)
	=
	\frac{1}{2\pi}\int_0^\pi
	\bigl(f(\theta-t)+f(\theta+t)\bigr)F_N(t)\dd t.
	\]
	同时
	\[
	L
	=
	\frac{1}{2\pi}\int_0^\pi
	(\alpha+\beta)F_N(t)\dd t.
	\]
	所以
	\[
	\begin{aligned}
		\sigma_N(f)(\theta)-L
		&=\frac{1}{2\pi}\int_0^\pi
		\bigl[
		f(\theta+t)-\alpha
		+f(\theta-t)-\beta
		\bigr]F_N(t)\dd t.
	\end{aligned}
	\]
	
	给定 $\varepsilon>0$,取与第 (a) 问相同的 $\delta>0$,使得
	当 $0<t<\delta$ 时,
	\[
	|f(\theta+t)-\alpha|<\varepsilon,
	\qquad
	|f(\theta-t)-\beta|<\varepsilon.
	\]
	
	于是
	\[\frac{1}{2\pi}\int_0^\delta
		\left|
		f(\theta+t)-\alpha
		+f(\theta-t)-\beta
		\right|F_N(t)\dd t
		\leq
		\frac{2\varepsilon}{2\pi}
		\int_0^\delta F_N(t)\dd t\leq
		\frac{2\varepsilon}{2\pi}
		\int_0^\pi F_N(t)\dd t
		=\varepsilon.
	\]
	
	在 $\delta\leq t\leq\pi$ 上
	\[
	F_N(t)
	\leq
	\frac{1}{N\sin^2(\delta/2)}.
	\]
	因此
	\[
	\sup_{\delta\leq t\leq\pi}F_N(t)
	\leq
	\frac{1}{N\sin^2(\delta/2)}
	\to0,
	\qquad N\to\infty.
	\]
	
	利用第 (a) 问中定义的常数 $C_\delta$,有
	\[\frac{1}{2\pi}\int_\delta^\pi
		\left|
		f(\theta+t)-\alpha
		+f(\theta-t)-\beta
		\right|F_N(t)\dd t\leq
		\frac{C_\delta}{2\pi}
		\sup_{\delta\leq t\leq\pi}F_N(t)
		\to0.
	\]
	
	所以
	\[
	\limsup_{N\to\infty}
	|\sigma_N(f)(\theta)-L|
	\leq\varepsilon.
	\]
	由于 $\varepsilon>0$ 任意,得到
	\[
	\lim_{N\to\infty}\sigma_N(f)(\theta)
	=L
	=\frac{f(\theta^+)+f(\theta^-)}{2}.
	\]

\end{proof}
对上面的Tauber定理，Littlewood 给出了 一个细化版本
\begin{proposition}
	\begin{enumerate}
		\item[(a)] 若级数
		\[
		\sum_{n=1}^{\infty}c_n
		\]
		的 Abel 和为 $s$，且
		\[
		c_n=O\left(\frac{1}{n}\right),
		\]
		则
		\[
		\sum_{n=1}^{\infty}c_n
		\]
		在通常意义下收敛于 $s$。
		
		\item[(b)] 从而，若级数
		\[
		\sum_{n=1}^{\infty}c_n
		\]
		的 Ces\`aro 和为 $s$，且
		\[
		c_n=O\left(\frac{1}{n}\right),
		\]
		则
		\[
		\sum_{n=1}^{\infty}c_n
		\]
		在通常意义下收敛于 $s$。
	\end{enumerate}
\end{proposition}
证明极其复杂，这里略过

这些结果可以应用到 Fourier 级数上。若可积函数 $f$
满足
\[
\widehat{f}(\nu)
=O\left(\frac{1}{|\nu|}\right),
\qquad |\nu|\to\infty,
\]
则有：

\begin{enumerate}
	\item[(i)] 若 $f$ 在 $\theta$ 处连续，则当 $N\to\infty$ 时，
	\[
	S_N(f)(\theta)\to f(\theta).
	\]
	
	\item[(ii)] 若 $f$ 在 $\theta$ 处是跃点，则当 $N\to\infty$ 时，
	\[
	S_N(f)(\theta)
	\to
	\frac{f(\theta^+)+f(\theta^-)}{2}.
	\]
	
	\item[(iii)] 若 $f$ 在 $[-\pi,\pi]$ 上连续，则
	\[
	S_N(f)\rightrightarrows f.
	\]
\end{enumerate}
%
